\section{Weekly Summary}

\begin{tcolorbox}[colback=yellow!10,colframe=black!50,title=AI-Generated Content]
\noindent The summary below was written by Claude (Anthropic) from that week's regression outputs and series descriptions. It has not been reviewed or fact-checked by a human, and should be read as an automated, informal recap -- not analysis.
\end{tcolorbox}

Welcome back to the only weekly column brave enough to shove two economic time series into a blender and act shocked when a smoothie comes out. As always, our beloved dartboard selected three pairs of variables that have absolutely no business being in the same spreadsheet, and as always, we ran the numbers anyway. Standard disclaimer, delivered in my most soothing voice: these are randomly paired series, nobody is claiming anyone causes anyone, and if you build a portfolio around this newsletter you deserve what happens next.

Monday gave us the Fed's holdings of medium-term Treasury securities versus a discontinued consumer price index for services in the South, Size Class D -- a series so obscure and so retired it hasn't posted since 2017. The raw regression came out gasping "significant!" with an R-squared of 0.77, and I was briefly tempted to write a think piece. Then we detrended it, the R-squared slumped to 0.37, and the slope politely flipped sign from positive to negative, which is the statistical equivalent of a witness changing their entire story on the stand. Translation: most of that "relationship" was just two things drifting through time together like strangers on the same escalator. A trend mirage, gently evaporating on contact. Deeply moving. Not real.

Tuesday, however, is where I have to reluctantly stop being sarcastic for approximately one sentence. We paired "Withdrawals from Income of Quasi-Corporations" (a phrase that sounds like a threat) against depreciation of nonresidential structures, and the raw fit was a frankly obnoxious R-squared of 0.98. My eyes rolled so hard I strained something -- surely another shared-trend fraud. But no: after detrending, it still posted an R-squared of 0.87 with a p-value that requires scientific notation to even write down, and the slope barely budged from 0.42 to 0.40. That is the rare, genuinely interesting case where the relationship survives having its trend surgically removed. Is there a real link? Plausibly -- both are dollar-denominated measures of nonfinancial business activity, so they may actually be tracking the same underlying economic thing rather than just aging in parallel. I refuse to be excited, but I am, quietly, in a corner.

And then Wednesday arrived to restore balance to the universe by regressing the population of Nuckolls County, Nebraska against a medical price index for congenital anomalies, because sure, why not. The raw fit strutted in at 0.88 with a confident negative slope, suggesting that as one Nebraska county empties out, the cost of certain medical services climbs -- a correlation so gloriously unrelated it belongs in a museum. Detrended, the R-squared collapsed to 0.36, though I'll note the slope stubbornly stayed negative and significant-ish, so this one is at least a more dignified mirage than Monday's. Still: two things going opposite directions over decades is not a relationship, it's a coincidence with good posture. See you next week, when the dartboard once again humbles us all.
